\section{Interpretation, prior-art boundary, and limitations}
\label{sec:interpretation}

Equation~\eqref{eq:headline} is best interpreted as a passive resource statement.  Resonances can redistribute gravitational response in frequency, internal hybridization can create bright and dark combinations, and matching choices can determine how much of a particular resonance is accessible from a selected local port.  What they cannot do, within the present class, is create additional cumulative gravitational coupling trace beyond the endpoint second-moment resource.  Free-space propagation then imposes the separate compact TT channel ceiling.

This perspective makes the minimum in Eq.~\eqref{eq:headline} physically important.  Either endpoint is a cut: a large source resource cannot compensate for a receiver whose total gravitational coupling trace is smaller, and vice versa.  The final result therefore eliminates both endpoint matching details rather than optimizing one resonator architecture.

The historical boundary is substantial.  Resonant-mass absorption cross sections and integrated response are old results \cite{PaikWagoner1976,Aguiar2010}; arbitrary-body multimode gravitational response is old \cite{Lobo1995}; material-response dispersion and sum-rule methods are old \cite{Srivastava2003}; and complete gravitational generator--receiver Hertz analyses are old \cite{Rudenko2003}.  Generic source--receiver channel decompositions and response-plus-propagation bounds are also established \cite{Miller2000,Molesky2020}.  The present work therefore makes no standalone novelty claim for the $20/3$ tensor contraction, the $4/3$ endpoint resource, passive $H_2$ machinery, directivity, or multiple scattering.

The narrower claimed contribution is the cumulative two-ended closure: the selected-port spectral area is cut by the smaller gravitational trace, each trace is reduced to $I_2$, the compact TT channel contributes the $25/16$ propagation coefficient, and repeated passive returns between the same endpoints do not change the leading asymptotic coefficient.  We have not found an equivalent complete theorem in the inspected literature; that negative search is not a priority claim.

The theorem has deliberate exclusions.  It does not cover active gain, inversion, parametric pumping, externally powered feedback, extended phased apertures, added gravitational relays or external cavities, reactive near-field exchange, higher-multipole-dominated operation, relativistic or nonlinear matter, arbitrary unbounded PDE boundary ports without an admissibility analysis, genuinely non-Markov continua, or curved-background focusing.  It is also a narrowband carrier-envelope result.  A broad absolute-frequency theorem would have to retain the explicit frequency dependence of the matter and propagation resources rather than pulling them outside the integral.

Within these limits, Eq.~\eqref{eq:headline} provides a compact answer to a practical theoretical question: passive resonant complexity can reshape the gravitational transfer spectrum, but it cannot increase its integrated two-ended ceiling beyond a resource fixed by endpoint mass distribution and separated TT propagation.

\section{Conclusion}

For a separated compact passive gravitational link, passivity first reduces frequency-integrated coherent transfer to the smaller gravitational coupling trace of the two endpoints.  Mass-weighted quadrupole completeness bounds each trace by $(4G/3c^5)I_2\omega_0^4$ in a narrow carrier band, while normalized compact TT propagation contributes the leading factor $25/[16(k_0R)^2]$.  Their combination gives
\begin{equation}
\Gamma_{\rm coh}
\lesssim
\frac{25G\omega_0^2}{12c^3R^2}
\min(I_{2,A},I_{2,B}).
\end{equation}
The same endpoint cut extends to countably infinite bounded-port Markov modal sectors, and repeated passive returns between the same compact endpoints do not alter the leading wave-zone coefficient.  The result is a specialist passive-resource theorem built from largely historical ingredients; its intended contribution is the explicit gravity-specific two-ended inertia closure.
